\documentclass[12pt,a4paper]{article}

\usepackage{fontspec}
\usepackage{polyglossia}
\setmainlanguage{russian}
\setotherlanguage{english}
\setmainfont{Times New Roman}
\setsansfont{Calibri}
\setmonofont{Courier New}

\usepackage{amsmath,amssymb,mathtools}
\usepackage{geometry}
\usepackage{booktabs}
\usepackage{hyperref}
\usepackage{enumitem}
\usepackage{physics}

\geometry{margin=2.5cm}

\title{Применение метода интегральных соотношений Баренблатта\\
к двухфазной black-oil модели нефти с растворённым газом}
\author{}
\date{}

\begin{document}

\maketitle

\section{Постановка задачи}

Рассматривается изотермическая радиальная фильтрация нефти с растворённым газом к добывающей скважине.
Воды в системе нет. Изначально пласт находится выше давления насыщения:
\begin{equation}
    p_i > p_b, \qquad S_g = 0, \qquad R_s = R_{si}.
\end{equation}

При снижении давления ниже давления насыщения около скважины начинается выделение газа в отдельную фазу:
\begin{equation}
    p < p_b, \qquad S_g > 0.
\end{equation}

Фазовые скорости задаются законом Дарси с относительными фазовыми проницаемостями:
\begin{equation}
    u_o = - \frac{k k_{ro}(S_g)}{\mu_o(p)} \pdv{p}{r},
\end{equation}
\begin{equation}
    u_g = - \frac{k k_{rg}(S_g)}{\mu_g(p)} \pdv{p}{r}.
\end{equation}

Здесь $u_o$ и $u_g$ --- объёмные скорости нефтяной и газовой фаз в пластовых условиях, $k$ --- абсолютная проницаемость, $k_{ro}$ и $k_{rg}$ --- относительные фазовые проницаемости, $\mu_o$ и $\mu_g$ --- вязкости фаз.

Предполагается отсутствие капиллярного давления и гравитации, поэтому давления в фазах одинаковы:
\begin{equation}
    p_o = p_g = p.
\end{equation}

Так как воды нет,
\begin{equation}
    S_o + S_g = 1.
\end{equation}

\section{Структура областей фильтрации}

Метод интегральных соотношений удобно применять через разбиение пласта на несколько зон с движущимися границами.

После начала разгазирования структура задачи имеет вид:
\begin{equation}
    r_w < r < r_b(t),
\end{equation}
где имеется двухфазная насыщенная зона:
\begin{equation}
    p < p_b, \qquad S_g > 0,
\end{equation}

\begin{equation}
    r_b(t) < r < R(t),
\end{equation}
где имеется однофазная возмущённая зона:
\begin{equation}
    p_b < p < p_i, \qquad S_g = 0,
\end{equation}

и область
\begin{equation}
    r > R(t),
\end{equation}
которая считается невозмущённой:
\begin{equation}
    p = p_i, \qquad S_g = 0.
\end{equation}

Здесь:
\begin{itemize}[noitemsep]
    \item $r_w$ --- радиус скважины;
    \item $r_b(t)$ --- радиус фронта разгазирования, где $p=p_b$;
    \item $R(t)$ --- внешний радиус возмущённой области, где $p=p_i$ и поток равен нулю.
\end{itemize}

Схематически:
\begin{equation}
    r_w
    \; | \;
    \text{нефть + свободный газ}
    \; | \;
    r_b(t)
    \; | \;
    \text{однофазная нефть с растворённым газом}
    \; | \;
    R(t)
    \; | \;
    \text{невозмущённый пласт}.
\end{equation}

\section{Black-oil уравнения без воды}

В радиальной постановке баланс нефтяной компоненты записывается как
\begin{equation}
    \pdv{}{t}
    \left(
        \phi \frac{S_o}{B_o}
    \right)
    +
    \frac{1}{r}\pdv{}{r}
    \left(
        r \frac{u_o}{B_o}
    \right)
    =0.
\end{equation}

Баланс газовой компоненты:
\begin{equation}
    \pdv{}{t}
    \left[
        \phi
        \left(
            R_s \frac{S_o}{B_o}
            +
            \frac{S_g}{B_g}
        \right)
    \right]
    +
    \frac{1}{r}\pdv{}{r}
    \left[
        r
        \left(
            R_s \frac{u_o}{B_o}
            +
            \frac{u_g}{B_g}
        \right)
    \right]
    =0.
\end{equation}

Здесь $B_o(p)$, $B_g(p)$, $R_s(p)$ --- PVT-функции black-oil модели.

В однофазной области $p>p_b$:
\begin{equation}
    S_g = 0, \qquad S_o = 1, \qquad R_s = R_{si}.
\end{equation}

В насыщенной области $p \leq p_b$:
\begin{equation}
    R_s = R_s(p), \qquad S_g > 0.
\end{equation}

Таким образом, в насыщенной зоне неизвестны две функции:
\begin{equation}
    p(r,t), \qquad S_g(r,t),
\end{equation}
а в однофазной зоне --- только давление $p(r,t)$.

\section{Аппроксимация решения в насыщенной зоне}

В двухфазной зоне $r_w<r<r_b(t)$ введём логарифмическую безразмерную координату:
\begin{equation}
    \xi = \frac{\ln(r/r_w)}{\ln(r_b/r_w)},
    \qquad
    0 \leq \xi \leq 1.
\end{equation}

Приближённый профиль давления можно записать в виде
\begin{equation}
    p_1(r,t)
    =
    p_w(t)
    +
    \left[p_b - p_w(t)\right] f_p(\xi).
\end{equation}

Простейший выбор:
\begin{equation}
    f_p(\xi)=\xi.
\end{equation}

Тогда выполняются условия
\begin{equation}
    p_1(r_w,t)=p_w(t),
    \qquad
    p_1(r_b,t)=p_b.
\end{equation}

Профиль газонасыщенности зададим как
\begin{equation}
    S_g(r,t)
    =
    S_{gw}(t) f_s(\xi).
\end{equation}

Например,
\begin{equation}
    f_s(\xi)=(1-\xi)^m,
    \qquad m \geq 1.
\end{equation}

Тогда
\begin{equation}
    S_g(r_w,t)=S_{gw}(t),
    \qquad
    S_g(r_b,t)=0.
\end{equation}

Параметр $m$ задаёт форму профиля газонасыщенности. В первом приближении можно брать $m=1$ или $m=2$.

\section{Аппроксимация решения в однофазной зоне}

В однофазной возмущённой зоне $r_b(t)<r<R(t)$ введём координату
\begin{equation}
    \eta = \frac{\ln(r/r_b)}{\ln(R/r_b)},
    \qquad
    0 \leq \eta \leq 1.
\end{equation}

Профиль давления зададим как
\begin{equation}
    p_2(r,t)
    =
    p_b + (p_i-p_b)g(\eta).
\end{equation}

Один из простейших вариантов:
\begin{equation}
    g(\eta)=2\eta-\eta^2.
\end{equation}

Тогда выполняются условия
\begin{equation}
    p_2(r_b,t)=p_b,
    \qquad
    p_2(R,t)=p_i,
\end{equation}
и
\begin{equation}
    \left.\pdv{p_2}{r}\right|_{r=R}=0.
\end{equation}

Последнее условие означает отсутствие потока на внешней границе возмущённой области.

\section{Интегральные соотношения в насыщенной зоне}

В двухфазной зоне введём невязку уравнения баланса нефтяной компоненты:
\begin{equation}
    \mathcal{R}_o
    =
    \pdv{}{t}
    \left(
        \phi \frac{1-S_g}{B_o}
    \right)
    +
    \frac{1}{r}\pdv{}{r}
    \left(
        r \frac{u_o}{B_o}
    \right).
\end{equation}

Невязка уравнения баланса газовой компоненты:
\begin{equation}
    \mathcal{R}_g
    =
    \pdv{}{t}
    \left[
        \phi
        \left(
            R_s \frac{1-S_g}{B_o}
            +
            \frac{S_g}{B_g}
        \right)
    \right]
    +
    \frac{1}{r}\pdv{}{r}
    \left[
        r
        \left(
            R_s \frac{u_o}{B_o}
            +
            \frac{u_g}{B_g}
        \right)
    \right].
\end{equation}

В первом приближении требуем выполнения интегральных балансов:
\begin{equation}
    \int_{r_w}^{r_b(t)} r\mathcal{R}_o\,dr = 0,
\end{equation}
\begin{equation}
    \int_{r_w}^{r_b(t)} r\mathcal{R}_g\,dr = 0.
\end{equation}

В более высоком приближении можно добавить весовые функции. Например:
\begin{equation}
    \int_{r_w}^{r_b(t)} r\xi\mathcal{R}_o\,dr = 0,
\end{equation}
\begin{equation}
    \int_{r_w}^{r_b(t)} r\xi\mathcal{R}_g\,dr = 0.
\end{equation}

В общем виде:
\begin{equation}
    \int_{r_w}^{r_b(t)} r\psi_j(\xi)\mathcal{R}_o\,dr = 0,
\end{equation}
\begin{equation}
    \int_{r_w}^{r_b(t)} r\psi_j(\xi)\mathcal{R}_g\,dr = 0,
\end{equation}
где $\psi_j(\xi)$ --- выбранные весовые функции, например $1,\xi,\xi^2,\ldots$.

\section{Интегральное соотношение в однофазной зоне}

В однофазной зоне $r_b(t)<r<R(t)$ имеем $S_g=0$, $S_o=1$, $R_s=R_{si}$.

Невязку можно записать по нефтяной компоненте:
\begin{equation}
    \mathcal{R}_{o,2}
    =
    \pdv{}{t}
    \left(
        \phi \frac{1}{B_o}
    \right)
    +
    \frac{1}{r}\pdv{}{r}
    \left(
        r \frac{u_o}{B_o}
    \right).
\end{equation}

Тогда первое интегральное соотношение:
\begin{equation}
    \int_{r_b(t)}^{R(t)} r\mathcal{R}_{o,2}\,dr = 0.
\end{equation}

При необходимости можно добавить более высокие соотношения:
\begin{equation}
    \int_{r_b(t)}^{R(t)} r\theta_j(\eta)\mathcal{R}_{o,2}\,dr = 0,
\end{equation}
где $\theta_j(\eta)=1,\eta,\eta^2,\ldots$.

\section{Скважинное условие}

Если задан дебит нефти в поверхностных условиях, то условие на скважине можно записать как
\begin{equation}
    q_o^{sc}
    =
    -2\pi h r_w
    \left.
    \frac{u_o}{B_o}
    \right|_{r=r_w}.
\end{equation}

Если задано забойное давление, то
\begin{equation}
    p(r_w,t)=p_w.
\end{equation}

Газовый дебит в поверхностных условиях:
\begin{equation}
    q_g^{sc}
    =
    -2\pi h r_w
    \left.
    \left(
        R_s \frac{u_o}{B_o}
        +
        \frac{u_g}{B_g}
    \right)
    \right|_{r=r_w}.
\end{equation}

Добываемый газовый фактор:
\begin{equation}
    GOR(t)=\frac{q_g^{sc}}{q_o^{sc}}.
\end{equation}

\section{Неизвестные функции и система ОДУ/ДАУ}

Для простейшего варианта с заданным дебитом нефти неизвестными могут быть
\begin{equation}
    r_b(t), \qquad R(t), \qquad p_w(t), \qquad S_{gw}(t).
\end{equation}

Для варианта с заданным забойным давлением $p_w$ неизвестными будут, например,
\begin{equation}
    r_b(t), \qquad R(t), \qquad S_{gw}(t),
\end{equation}
а дебиты $q_o^{sc}(t)$ и $q_g^{sc}(t)$ будут вычисляться из градиентов давления на скважине.

После подстановки аппроксимаций $p_1(r,t)$, $S_g(r,t)$ и $p_2(r,t)$ в невязки $\mathcal{R}_o$, $\mathcal{R}_g$ и $\mathcal{R}_{o,2}$ интегральные соотношения дают систему ОДУ/ДАУ относительно параметров профиля:
\begin{equation}
    \mathbf{F}
    \left(
        r_b,
        R,
        p_w,
        S_{gw},
        \dot r_b,
        \dot R,
        \dot p_w,
        \dot S_{gw},
        t
    \right)
    =0.
\end{equation}

В общем случае эта система решается численно.


\section{Адаптационные параметры и параметры формы}

В предложенной постановке метод интегральных соотношений содержит не только физические параметры пласта и флюидов, но и параметры, связанные с выбранной формой аппроксимирующих профилей. Такие параметры удобно называть адаптационными параметрами или параметрами формы.

Их появление связано с тем, что вместо полного решения системы дифференциальных уравнений black-oil модели заранее выбирается конечномерный вид функций $p(r,t)$ и $S_g(r,t)$. Поэтому часть информации о реальной форме решения переносится в выбранные функции $f_p(\xi)$, $f_s(\xi)$ и $g(\eta)$.

\subsection{Физические параметры}

К физическим параметрам относятся величины, которые должны задаваться по данным пласта, PVT-исследований, относительных фазовых проницаемостей или по результатам гидродинамической интерпретации:
\begin{equation}
    k, \qquad h, \qquad \phi, \qquad c_r,
\end{equation}
\begin{equation}
    B_o(p), \qquad B_g(p), \qquad R_s(p), \qquad \mu_o(p), \qquad \mu_g(p),
\end{equation}
\begin{equation}
    k_{ro}(S_g), \qquad k_{rg}(S_g), \qquad S_{gc}, \qquad S_{or}.
\end{equation}

Эти параметры не являются параметрами формы метода. Их желательно не использовать для подгонки самой интегральной модели, если для них есть независимые лабораторные или промысловые оценки.

\subsection{Динамические неизвестные}

Величины
\begin{equation}
    R(t), \qquad r_b(t), \qquad p_w(t), \qquad S_{gw}(t)
\end{equation}
не являются адаптационными параметрами. Это динамические неизвестные, которые должны быть найдены из интегральных соотношений, граничных условий и скважинного условия.

Если учитывается критическая газонасыщенность, то к динамическим неизвестным добавляется фронт подвижного газа:
\begin{equation}
    r_m(t).
\end{equation}

Таким образом, параметры типа $R(t)$, $r_b(t)$ и $S_{gw}(t)$ являются результатом расчёта, а не свободными настроечными коэффициентами.

\subsection{Параметры формы}

К адаптационным параметрам относятся показатели и коэффициенты, задающие форму профилей давления и насыщенности. Например, если газонасыщенность в насыщенной зоне задаётся как
\begin{equation}
    S_g(r,t)=S_{gw}(t)
    \left[
        1-
        \frac{\ln(r/r_w)}{\ln(r_b/r_w)}
    \right]^m,
\end{equation}
то показатель $m$ является параметром формы. Он не является прямым физическим параметром пласта. Его смысл --- степень локализации свободного газа около скважины.

При $m=1$ профиль газонасыщенности более пологий:
\begin{equation}
    S_g \sim 1-\xi.
\end{equation}

При $m=2$ свободный газ сильнее сконцентрирован около скважины:
\begin{equation}
    S_g \sim (1-\xi)^2.
\end{equation}

При больших значениях $m$ газонасыщенность ещё сильнее локализуется в прискважинной области.

Аналогичная ситуация возникает для профиля давления. Например, вместо линейного профиля
\begin{equation}
    f_p(\xi)=\xi
\end{equation}
можно использовать более общий вид
\begin{equation}
    f_p(\xi)=1-(1-\xi)^n,
\end{equation}
где $n$ также является параметром формы.

Тогда давление в насыщенной зоне можно записать как
\begin{equation}
    p_1(r,t)
    =
    p_w(t)+\left[p_b-p_w(t)\right]
    \left[1-(1-\xi)^n\right].
\end{equation}

Для однофазной зоны аналогичным образом можно ввести семейство функций
\begin{equation}
    p_2(r,t)=p_b+(p_i-p_b)g(\eta),
\end{equation}
например
\begin{equation}
    g(\eta)=1-(1-\eta)^\alpha,
\end{equation}
или другую функцию, удовлетворяющую условиям
\begin{equation}
    g(0)=0, \qquad g(1)=1, \qquad \left.\dv{g}{\eta}\right|_{\eta=1}=0
\end{equation}
при необходимости нулевого потока на внешней границе возмущённой области.

Таким образом, типичные адаптационные параметры можно обозначить как
\begin{equation}
    m, \qquad n, \qquad \alpha, \qquad \beta, \ldots
\end{equation}
Они определяют не физические свойства пласта, а выбранную форму аппроксимации решения.

\subsection{Выбор параметров формы}

Для первого инженерного варианта желательно использовать минимальное число адаптационных параметров. Например, можно зафиксировать профили давления и оставить только один параметр формы для газонасыщенности:
\begin{equation}
    p_1=p_w+(p_b-p_w)\xi,
\end{equation}
\begin{equation}
    S_g=S_{gw}(1-\xi)^m,
\end{equation}
\begin{equation}
    p_2=p_b+(p_i-p_b)(2\eta-\eta^2).
\end{equation}

Тогда единственным настраиваемым параметром остаётся $m$.

Возможная интерпретация:
\begin{equation}
    m=1
    \quad \text{--- относительно равномерное распределение свободного газа},
\end{equation}
\begin{equation}
    m=2
    \quad \text{--- умеренная концентрация газа около скважины},
\end{equation}
\begin{equation}
    m>2
    \quad \text{--- сильная локализация газа в прискважинной зоне}.
\end{equation}

Подбор параметров формы следует выполнять не по каждой отдельной точке времени, а по интегральным характеристикам решения:
\begin{equation}
    q_o(t), \qquad q_g(t), \qquad GOR(t), \qquad r_b(t), \qquad p_w(t).
\end{equation}

В качестве эталона можно использовать одномерную радиальную black-oil модель, численный гидродинамический симулятор или предельные аналитические решения для упрощённых случаев. Если параметры формы подгоняются отдельно для каждого момента времени, модель теряет предсказательность и превращается в чисто аппроксимационную процедуру.

\subsection{Вариант без явных адаптационных показателей}

В более строгом варианте метода можно отказаться от фиксированных показателей $m$ и $n$. Вместо этого профили давления и газонасыщенности задаются полиномами с неизвестными коэффициентами:
\begin{equation}
    p(r,t)=a_0(t)+a_1(t)\xi+a_2(t)\xi^2+\ldots+a_N(t)\xi^N,
\end{equation}
\begin{equation}
    S_g(r,t)=b_0(t)+b_1(t)\xi+b_2(t)\xi^2+\ldots+b_M(t)\xi^M.
\end{equation}

Коэффициенты $a_i(t)$ и $b_i(t)$ определяются из граничных условий и расширенного набора интегральных соотношений:
\begin{equation}
    \int r\mathcal{R}\,dr=0,
\end{equation}
\begin{equation}
    \int r\xi\mathcal{R}\,dr=0,
\end{equation}
\begin{equation}
    \int r\xi^2\mathcal{R}\,dr=0,
\end{equation}
и так далее.

В этом случае форма профиля не задаётся одним адаптационным показателем, а определяется самой системой уравнений. Такой подход ближе к классической идее метода интегральных соотношений, но приводит к более сложной системе ОДУ/ДАУ.

\subsection{Практический вывод}

В упрощённой МИС-модели для двухфазной black-oil задачи адаптационные параметры действительно есть. В первую очередь это параметры формы профилей давления и газонасыщенности. Наиболее важным из них в первом приближении является показатель $m$ в профиле $S_g(r,t)$, поскольку именно он управляет локализацией свободного газа около скважины и существенно влияет на прогноз $GOR(t)$.

\section{Учёт критической газонасыщенности}

Если свободный газ становится подвижным только после достижения критической газонасыщенности $S_{gc}$, то
\begin{equation}
    k_{rg}=0, \qquad S_g<S_{gc}.
\end{equation}

Тогда полезно ввести дополнительный фронт
\begin{equation}
    r_m(t): \qquad S_g(r_m,t)=S_{gc}.
\end{equation}

Структура зон становится следующей:
\begin{equation}
    r_w<r<r_m(t):
    \qquad
    \text{нефть + подвижный газ},
\end{equation}
\begin{equation}
    r_m(t)<r<r_b(t):
    \qquad
    \text{нефть + неподвижный свободный газ},
\end{equation}
\begin{equation}
    r_b(t)<r<R(t):
    \qquad
    \text{однофазная нефть с растворённым газом},
\end{equation}
\begin{equation}
    r>R(t):
    \qquad
    \text{невозмущённый пласт}.
\end{equation}

В первом инженерном приближении можно принять $S_{gc}=0$. Но для реалистичного прогноза газового фактора $GOR(t)$ критическую газонасыщенность желательно учитывать.

\section{Практический алгоритм расчёта}

Расчёт можно организовать следующим образом:

\begin{enumerate}[label=\arabic*.]
    \item Задать PVT-функции:
    \begin{equation}
        B_o(p), \qquad B_g(p), \qquad \mu_o(p), \qquad \mu_g(p), \qquad R_s(p).
    \end{equation}

    \item Задать относительные фазовые проницаемости:
    \begin{equation}
        k_{ro}(S_g), \qquad k_{rg}(S_g).
    \end{equation}

    \item До достижения $p_w=p_b$ решать однофазную задачу методом интегральных соотношений.

    \item После появления области $p<p_b$ ввести фронт разгазирования $r_b(t)$.

    \item Задать приближённые профили:
    \begin{equation}
        p_1(r,t), \qquad S_g(r,t), \qquad p_2(r,t).
    \end{equation}

    \item Подставить профили в компонентные black-oil балансы.

    \item Составить интегральные соотношения:
    \begin{equation}
        \int r\mathcal{R}_o\,dr=0,
        \qquad
        \int r\mathcal{R}_g\,dr=0.
    \end{equation}

    \item Получить систему ОДУ/ДАУ для параметров профиля:
    \begin{equation}
        r_b(t), \qquad R(t), \qquad p_w(t), \qquad S_{gw}(t).
    \end{equation}

    \item Численно интегрировать систему по времени.

    \item Рассчитывать добычные показатели:
    \begin{equation}
        q_o(t), \qquad q_g(t), \qquad GOR(t), \qquad r_b(t), \qquad R(t).
    \end{equation}
\end{enumerate}

\section{Физическая интерпретация}

Пока давление во всём пласте выше давления насыщения, течение фактически однофазное. Нефть содержит растворённый газ, но свободная газовая фаза отсутствует. Дебит определяется однофазной пьезопроводностью и PVT-свойствами нефти выше давления насыщения.

После снижения давления около скважины ниже $p_b$ начинается выделение газа:
\begin{equation}
    R_s(p) < R_{si}.
\end{equation}

Часть газовой компоненты переходит из растворённого состояния в свободную фазу. Это приводит к появлению газонасыщенности $S_g$ и изменению фазовых подвижностей:
\begin{equation}
    \lambda_o = \frac{k_{ro}(S_g)}{\mu_o(p)},
    \qquad
    \lambda_g = \frac{k_{rg}(S_g)}{\mu_g(p)}.
\end{equation}

Обычно рост $S_g$ снижает нефтяную относительную проницаемость $k_{ro}$ и увеличивает газовую относительную проницаемость $k_{rg}$. Поэтому ожидаемые эффекты:
\begin{equation}
    q_o(t) \downarrow,
    \qquad
    GOR(t) \uparrow.
\end{equation}

Если есть критическая газонасыщенность, то свободный газ сначала может быть неподвижным. В этом случае добываемый газовый фактор сначала может оставаться близким к растворённому, а затем резко увеличиваться после появления подвижного газа у скважины.

\section{Вывод}

Метод интегральных соотношений можно применить к двухфазной black-oil задаче, но он уже не сводится к одному уравнению пьезопроводности. Необходимо рассматривать систему компонентных балансов и несколько движущихся границ:
\begin{equation}
    R(t) \quad \text{--- внешний фронт возмущения давления},
\end{equation}
\begin{equation}
    r_b(t) \quad \text{--- фронт разгазирования},
\end{equation}
\begin{equation}
    r_m(t) \quad \text{--- фронт подвижного газа, если учитывается } S_{gc}.
\end{equation}

В простейшем варианте модель имеет неизвестные
\begin{equation}
    R(t), \qquad r_b(t), \qquad p_w(t), \qquad S_{gw}(t),
\end{equation}
а интегральные уравнения получаются из балансов нефтяной и газовой компонент.

Такой подход не заменяет полноценный численный black-oil симулятор, но может быть полезен как компактная полуаналитическая модель для инженерной оценки радиуса разгазирования, изменения дебита нефти и роста газового фактора.

\end{document}
